% =====================================================================
% si_n3.tex -- Supporting Information: N3 preregistration, amendments and
% disclosures. P = tasks/n3_protocol.md. The disclosure list reproduces
% P Section 1.4 (a)-(h) with light copy-editing only. Requires booktabs.
% Citation placeholders use \CITEME{...}; si.tex must define it as main.tex
% does.
% =====================================================================
\section{Preregistration, amendments and disclosures}\label{si:prereg}

The N3 protocol (\texttt{tasks/n3\_protocol.md}, version 3) was committed and
pushed to the public repository at github.com/fronkt/symmc-flow on 2026-09-28 at
01:59 EDT (commit 8f7e494), before any N3 arm was trained, sampled or scored.
The push is the third-party timestamp. Version 1 had been reviewed from five
adversarial viewpoints (a MolCrystalFlow author, a statistician, a referee looking
for weak baselines, an implementing engineer and an ACS Omega referee), which
raised 48 issues; every blocker and major issue was re-checked against code and
data. Version 2 adopted the checked fixes, and a second review (38 issues) led to
version 3. A third pass on version 3, limited to blockers and majors, found four
major issues, whose checked fixes are part of the frozen text. After the freeze the
protocol changed only through dated amendments appended to the same file (A1, A2,
A2a, A3, A4 and the step-5b clarifications), each stating which N3 outputs existed
when it was made. The step-6 selection freeze records the SHA-256 of the protocol
file, and the scoring code refuses a freeze whose protocol has changed, so the
three later execution clarifications were kept in separate dated files.
Table~\ref{tab:si_timeline} lists the records in order.
% src: tasks/n3_protocol.md:1-17,968; tasks/n3_step7_clarification_2026-11-08.md:3-5; git log (commit dates)

\begin{table}
\caption{Preregistration and execution records, in commit order. Times are local
(EDT before 2026-11-01, EST after).}
\label{tab:si_timeline}
\small
\begin{tabular}{@{}p{5.2cm}p{2.6cm}p{1.5cm}p{6.0cm}@{}}
\toprule
record & committed & commit & state when made \\
\midrule
protocol version 3 frozen & 2026-09-28 01:59 & 8f7e494 & no N3 arm trained, sampled or scored \\
A1: SEL and TEST-B frozen & 2026-09-28 02:22 & 05aa3c9 & no arm output on SEL, TEST-B or DEV-TEST \\
A2: implementation clarifications & 2026-09-28 07:52 & 74f1921 & as A1; CPU smoke runs on VAL and TRAIN only \\
A2a: deployment packaging & 2026-10-11 01:06 & 783a1ab & as A1 \\
A3: step-4 settings and approval & 2026-10-19 15:16 & a25a444 & as A1; no VALSEL truth minimised \\
H2 form published (EXACT) & 2026-10-21 10:24 & 1aff111 & before any candidate was minimised under FF* \\
A4: Gate C corrections & 2026-11-02 16:32 & b50fb74 & as A1; no MCF arm trained \\
step-5b clarifications & 2026-11-06 17:22 & 159876d & before any VALSEL candidate score \\
step-6 selection freeze & 2026-11-08 21:54 & 772a36e & before any TEST-B or DEV-TEST input on the cluster \\
step-7 clarification & 2026-11-08 22:20 & 8e45e1a & as the freeze \\
step-7 code fix & 2026-11-09 13:19 & 412555f & as the freeze \\
step-9 clarification & 2026-11-15 19:32 & f45d80a & before any step-9 (scoring) command \\
results (\texttt{n3\_results.json}) & 2026-11-17 15:36 & f1c713a & before interpretation \\
\bottomrule
\end{tabular}
\end{table}
% src: git show -s for each commit; tasks/n3_protocol.md:1085-1087,1131,1328,1363,1425,1520; tasks/n3_step7_clarification_2026-11-08.md:1; tasks/n3_step7_codefix_2026-11-09.md:1; tasks/n3_step9_clarification_2026-11-15.md:1; tasks/n3_step4_record_2026-10-11.md:1043-1045

\subsection{Disclosures}\label{si:disclosures}

The protocol requires the following disclosures (its Section 1.4). DEV-TEST is
labelled ``used during development, not confirmatory'' wherever it is reported.
\begin{itemize}
\item[(a)] DEV-TEST was scored by the G2 baselines, the G2 evaluation, N1 and N2.
\item[(b)] The N1 idea of ranking draws before relaxing them came from a DEV-TEST
diagnostic: the relaxed-truth drift, median 14.4$^\circ$.
\item[(c)] The model of record for OURS was fixed after DEV-TEST had been seen. The
N2 retrain and the G2 checkpoints tied on VAL picked@1 (7.0 vs 7.0\% per-seed mean;
9.0 vs 9.0\% pooled). Both were then scored on DEV-TEST: G2 7.5\%, N2 5.0\%
(crystal-level sign-flip $p=0.023$). G2 was kept. Applied when the protocol was
written, and labelled post hoc, the rule ``incumbent unless VAL-better; tie
$\rightarrow$ VAL draw@1'' keeps G2 from VAL alone: draw@1 0.96 vs 0.67\%, higher
in every seed; any@16 13.3 vs 10.3\%. Applied per seed across the two runs, the
same rule would replace seed 1 with the N2 checkpoint (VAL 7 vs 4\%), and DEV-TEST
would be 6.5\% instead of 7.5\%. N2 is carried as a sensitivity family (protocol
Section 2.4).
\item[(d)] The DEV-TEST stratum rates of OURS were computed before the strata were
frozen: mixed-handedness 8.5\%, Sohncke 6.1\%, improper-achiral 3.5\% (per-seed
means). The frozen refined chirality definition moves the mixed-handedness rate of
OURS from 8.2\% (35/429, capped definition, 143 crystals) to 8.5\% (35/411, 137
crystals), because OURS picks 0 of 18 on the 6 reclassified crystals.
\item[(e)] During protocol review, DEV-TEST (and VAL) were also used for:
loose-tolerance rescoring of stored OURS draws; Haar floors at loose tolerance;
homochiral-approximant matches (1/131 strict vs 80/131 at stol 0.5, partial run);
MCF perfect-model orientation ceilings (conventional cells, keyed on $K$); a
perfect-model orbit-LJ check (4/28 $K=2$ crystals recovered); steric-press drift
checks; a VAL study of the iii-s screen (a start within 12.5$^\circ$ of the truth
is among the 64 lowest-LJ grid points in 91/98 VAL crystals); and X--H bond-length
statistics. None of these was used to choose anything in the protocol.
\item[(f)] The legacy det$+1$ pipeline parsed the TEST-B pool in July 2026, and the
old paper's scale runs trained on it. No legacy checkpoint is used (protocol
Section 1.2).
\item[(g)] Build step 1 (the rigid-molecule parse of the 4,500 pool CIFs) was
started on 2026-09-28, before the protocol was committed. It is input-side only
and produced no arm output. Its start time, command line and worker count are
recorded in amendment A1.
\item[(h)] H1--H3 are decided on TEST-B only. Everything on DEV-TEST is secondary.
\end{itemize}
% src: tasks/n3_protocol.md:164-198

\subsection{Amendments}\label{si:amendments}

Each amendment below states that it changes no hypothesis, comparator, set,
matcher or statistic; A3, A4 and the step-5b clarifications also state that they
change no selection menu, seed, endpoint, never-reduced item or
outcome-to-wording rule.
% src: tasks/n3_protocol.md:1135-1136,1349,1421,1518,1522-1523

\paragraph{A1 (2026-09-28): SEL and TEST-B frozen.} The SEL and TEST-B refcode lists
were committed with their SHA-256 hashes before any arm produced output on
them, together with the build counts given in the Methods and the hashes of the
parse cache, the UMA s-1p1 weights and element references, and the OMC25
starting-crystal list. For disclosure (g), it records that build step 1 ran from
00:14:57 to 01:28 EDT with five workers, ending before the protocol commit at
01:59:41 EDT. Before that commit, the expansion-check filter was checked on the
development corpus (VAL 100/100 and DEV-TEST 200/200 pass; TRAIN has the 9
predicted orbit failures, which stay in TRAIN).
% src: tasks/n3_protocol.md:1085-1127,1258-1260

\paragraph{A2 (2026-09-28): implementation clarifications.} Made while the only
outputs were CPU smoke runs on VAL and TRAIN and input-side gates, A2 filled gaps
and corrected facts found during implementation. The items a reader should know:
\begin{itemize}
\item The export gate (Gate F) passed on every set. The design power at
$n=1{,}000$ (4,000 simulated replicates, exact test; power at $\alpha=0.05$, at
$\alpha/3$ in parentheses) is 0.83 (0.71) at a 2 percentage-point difference and
0.97 (0.92) at 2.5 points for a three-seed comparator, 0.61 (0.45) and 0.82 (0.68)
for the single-run iii-s, and 0.74 (0.57) and 0.93 (0.85) for the H1
intersection-union test. The H2 comparator is therefore the least powered.
\item MCF's bond-perception step (RDKit\CITEME{RDKit} DetermineBonds inside MCF's
fix\_mol\_bonds) ran with a cap of $10^7$ iterations, because uncapped it does not
terminate on TRAIN crystal ZAYGOY. A capped molecule follows MCF's own fallback.
Only ZAYGOY (4 copies) hit the cap.
\item MCF's relative-rotation azimuth feature has a branch cut, and exact
C$_2$-conjugate copies, which every symmetric-prior draw has at $t=0$, sit on it.
The equivariance diagnostic (a) is therefore evaluated at a generic state, the
invariance of MCF's sampled distribution is approximate, and diagnostics (b)--(d)
are labelled approximate if the 50-step sampler deviation exceeds $10^{-3}$ on any
selected checkpoint. MCF's sampler runs only in float32.
\item MCF's own scorer rebuilds each lattice from its six parameters but keeps the
Cartesian coordinates in the original frame, so for cells not in pymatgen's
parameter orientation it reads a different structure from the one stored (VAL
68/100 and DEV-TEST 115/200 of our exports). Gate C-a (i) therefore compares MCF's
scorer with an emulation of that reading (\texttt{n3\_score.py --mcf-reading}).
Whether MCF's published numbers are affected is not claimed.
\item A draw with any non-finite lattice, centroid, rotation or coordinate is an
arm failure. A failing MCF inference call is retried, then sampled in chunks
of 25 crystals and then crystal by crystal with derived seeds; a crystal that
still fails is a miss, listed by refcode.
\item The MCF residual geodesic is computed in float64 with atan2. The ii-S
pull-back accuracy quoted in protocol Section 3.2 (``max 0.014$^\circ$'') was a
float32 floor; the true maxima are $8.1\times10^{-6}$ degrees (VAL, 100/100) and
$7.3\times10^{-6}$ degrees (DEV-TEST, 200/200).
\item MCF's cell standardization mirrors the cell when $\det Q=-1$ (VALSEL 142/400,
TEST-B 378/1,000). Every mapping uses the stored $Q$, and the matchers do not
distinguish enantiomorphs.
\item Coset keys unseen in TRAIN, which iv-P maps to a zeroed reserved label: SEL 2
copies in 1 crystal and TEST-B 27 copies in 6 crystals; none in VAL or DEV-TEST.
The legacy iv-A arm, which is not run, has 599 TRAIN crystals, not 598.
\item The FF* configuration is fixed once, at the step-4 timing gate, and is
identical for the truths, iii-s, i-E and the BASIN jobs on every set; a call that
returns a different resolved configuration is an environment error, never an FF*
failure. The finite-difference gradient check passes only if every relative error
is at most $10^{-4}$ in float64, on 5 VAL crystals at the true pose and one Haar
pose.
\item X--H normalisation moves each H bonded to C, N or O along its bond to exactly
the normalised length, lengthening or shortening it; other H atoms are unchanged.
\item Disclosure. Local CPU smoke runs minimised a few arm candidates under FF*
before any H2 form existed (iii-s starts on VAL crystals 26 and 78; i-E on the
iii-u draws of VAL crystal 26, and BASIN on its pick; in review, TRAIN crystal 57),
each with UMA s-1p1 in float32, at most 10 evaluations per start and no matching
of any kind. All of those outputs were deleted. They carry no information about the
H2 form, which is a fixed threshold on minimised truths.
\item OMC25's element set, taken as the union of elements over all its molecules,
is \{B, Br, C, Cl, F, H, I, N, O, P, S, Si\}. Crystals with other elements: DEV-TEST
1, VAL 2, SEL 5, TEST-B 16 (listed by refcode). FF* runs on them all.
\item For the stored N1 and N2 draws, the recorded steric energy (3 decimals) and
residual (5 decimals) are the selector inputs; recomputing the energy on the scored
structure changes no pick (all 100 near-tie crystal-seed pairs checked). A
deterministic energy failure gives a non-finite selector input, which ranks last.
\item For H1, the seed-sensitivity wording applies if either component's two-stage
interval contains zero. The standard deviation over seeds in Table 1 is the sample
value (ddof $=1$).
\item Table R column (c) uses 17 alignment starts (16 listed plus the best
automorphism rotation); the 16-start-only column is reported beside it.
\end{itemize}
% src: tasks/n3_protocol.md:1129-1324 (records 1138-1165; items 1, 3, 4, 7, 8, 9, 13, 18, 19, 20, 21, 27, 28, 29, 30, 31, 33, 34)

\paragraph{A2a (2026-10-11): deployment packaging and integrity.} The FF*
environment used a CUDA 13.0 PyTorch build, because no CUDA 12.8 wheel exists for
the pinned torch 2.13.0, so the GPU node needed a CUDA 13.0 or newer driver; the
MCF and SymMC environments kept CUDA 12.8. Transfers were staged so that TEST-B and
DEV-TEST inputs could reach the cluster only after the complete step-6 selection
was committed and pushed, and the selection and scoring tools check the hashes of
the selected checkpoints, configurations, truths and matcher libraries. No
scientific item changed.
% src: tasks/n3_protocol.md:1326-1359

\paragraph{A3 (2026-10-19): step-4 settings, approval and disclosures.} Made after
the timing gate (TRAIN and VAL crystals, nothing matched) and before any VALSEL
truth was minimised. The run moved from rented single GPU boxes to Purdue Anvil
(ACCESS allocation CHE260157, A100 nodes, one Apptainer image) with the author's
approval of a 700 GPU-h cap for steps 5--7, against a projection of 515--535 GPU-h;
no degrade item was taken, and on Anvil's persistent file system the step-8 ``destroy
the box'' became a cleanup script followed by a search of the home, scratch and
project file systems. FF* was fixed to fairchem's \texttt{batch} settings in
float32 (the other settings and float64 were measured and did not fit the cap),
the MCF speed-ups and iv-P's packed encoder were switched on after equivalence
checks, and a finite MCF draw with a non-finite residual stays valid for every
other selector and ranks last under RESID (counts are reported). The
finite-difference gradient check failed at step $h=10^{-4}$ rad on one of ten rows,
a clash pose whose error scaled as $h^2$ (central-difference truncation, not a
gradient defect), so the step was changed to $h=10^{-5}$ rad with the same pass
rule; the check had to pass again on the production GPU before any VALSEL truth was
minimised, and it did (commit 1aff111). A3 also discloses that one timing upload
was made in automatic mode on the author's prior approval instead of command by
command, and that the first timing box was abandoned and its measurements
discarded.
% src: tasks/n3_protocol.md:1361-1421,1425-1426; git 1aff111 (message: FD h=1e-5 passes on A100)

\paragraph{A4 (2026-11-02): Gate C corrections.} Gate C had run once
(2026-10-22) and recorded a failure, so no MCF arm had been trained. Gate C-a
had been scored against the wrong reading: the protocol (A2) gates check (i) on the
emulation of MCF's reading and check (ii) on MCF's own scorer, but the gate code had
never followed that text and gated both on our scorer's own reading. Re-scored from
the same prediction files, with nothing re-sampled, (i) agreed on every outcome and
(ii) fell within the $\pm5$ percentage-point window on both the TEST and VAL splits,
so the H1 harness clause is not triggered; that our scorer's own reading matched
more often than MCF's scorer on these predictions is reported as an observation,
not a claim. Gate C-b's GPU half failed under its original rule (the
patched-versus-unpatched difference was larger than an unpatched-versus-unpatched
difference for 10 of 24 nondeterministic tensors). Because both sides are single
draws of GPU nondeterminism, that rule fails an unchanged patch with high
probability (over 90\% on synthetic noise), so after the failure the author
replaced it, before re-running, with a noise envelope of one patched and 19
unpatched runs that fails only if the patched run is the strict maximum of the 20
(false-failure probability at most 5\%); the GPU half passed under the new rule,
and the CPU half, which had passed bitwise, remains the exact check. A4 also
relaxed the cluster prelude's test-path rule for Gate C's own run directories only.
% src: tasks/n3_protocol.md:1423-1518; results/n3/gateC.json (C-b.cuda: rule A4 item 2 (b), runs 20, pass true)

\paragraph{Step-5b clarifications (2026-11-06).} Recorded before any VALSEL
candidate was scored.
\begin{itemize}
\item The iv-P seed-0 run at learning rate $10^{-3}$ diverged in epoch 1, so its
validation-loss minimum is non-finite. A non-finite minimum ranks below every
finite one, as for every non-finite selector input. The choice was unchanged:
$10^{-4}$ (5.33186, against 5.33210 at $3\times10^{-4}$).
\item MCF's ``top-3 + last'' checkpoint candidates are read as a set. In all six
MCF runs the released checkpointer, which rewrites \texttt{last.ckpt} only when it
saves a top-3 checkpoint, left \texttt{last.ckpt} byte-identical to the
latest-epoch top-3 checkpoint, so no final-epoch (epoch 1,715) weights exist. Each
distinct checkpoint was sampled once, and the duplicated one carries the label
``last'' (Table~\ref{tab:si_mcf_ckpt}).
\item The iii-s VAL sanity row ran in step 5, where the projection booked it. It
chooses nothing.
\item The last MCF-A Stage-1 tie-break is MCF's logged validation loss; MCF-R
Stage B breaks ties on the fixed-draw VAL loss, as written.
\end{itemize}
% src: tasks/n3_protocol.md:1520-1541

\subsection{Execution clarifications after the selection freeze}\label{si:clarifications}

\paragraph{Step 7 (2026-11-08): stops on the cluster.} The protocol was written for
a single rented box: if the box dies during step 7, every partial TEST output is
discarded unscored and the full run list is repeated. On Anvil, step 7 ran as
several Slurm jobs. The clarification defines one job ending early (timeout, node
failure, preemption, out of memory) as not a box death: finished, checksummed runs
stay, resumable runs continue from their complete per-crystal records under a
configuration check, single-call runs restart from scratch with their committed
seeds, and every stop is logged. Loss or corruption of the box root, a failed
verification, a scratch purge, loss of the allocation or a stop forced by the cap
still mean the full discard and re-run. In the event, no step-7 job stopped and no
run was restarted.
% src: tasks/n3_step7_clarification_2026-11-08.md:1-35; tasks/n3_step4_record_2026-10-11.md:850-865

\paragraph{Step 7 (2026-11-09): code fix.} A smoke test on VAL and VALSEL inputs,
run before the upload, showed that every MCF-A prior-floor run would crash when its
outputs were bundled: an A3 change read an \texttt{info} field that the floor's
check function never returned. The fix makes that function return an empty
\texttt{info}; it changes no computation, run, seed or command. The fixed
\texttt{scripts/n3\_mcf\_run.py} (SHA-256 86dfe34d\ldots) was shipped with the
step-7 inputs and verified on the cluster, with a regression test added.
% src: tasks/n3_step7_codefix_2026-11-09.md:1-64

\paragraph{Step 9 (2026-11-15): scoring.} Step 9 ran as two Slurm jobs, one per
set, in Anvil's GPU partition because the allocation's CPU-only units were reserved
for other work; the GPU was not passed to the container, and no FF* evaluation ran
after step 8. Scoring used the environment VALSEL was scored in (PyTorch 2.11.0,
pymatgen 2026.5.4, SciPy 1.17.1, NumPy 2.3.5) and code byte-identical to the
freeze apart from the step-7 fix, and a driver ran only the frozen commands (92 per
set, with three permitted rewrites: worker count, interpreter path and freeze hash)
after checking that no scored output of the set existed. The clarification added
one aggregation that the protocol asks for (its Section 3.2, diagnostic (c)) but no
frozen command performed, a summary of the oracle-aligned MCF-R scores labelled an
oracle item; set a resume rule for stopped jobs (every step-9 command is
deterministic) and a rule that any code fix is recorded, pushed and disclosed; and
recorded in advance that the random floor's continuous-RMS block would read
``missing'', because the frozen list matches iii-u under that rule only at its LJ
picks. No run was added, and neither job stopped.
% src: tasks/n3_step9_clarification_2026-11-15.md:1-103; tasks/n3_step4_record_2026-10-11.md:1004-1020

Two step-9 deviations are disclosed. First, the results assembler
(\texttt{scripts/n3\_results.py}) stopped on its first run because it required the
H1--H3 families in the DEV-TEST report, which the frozen report code produces only
for TEST-B, by design. It was fixed (families required on TEST-B only), with two
tests added, before any number was read; only the reports' key structure was
inspected. The fix was committed together with the results (commit f1c713a), not
pushed before the re-run as the clarification's code-fix rule asks. Second, the
pullback of the step-9 outputs was run in automatic mode at the author's
instruction, rather than in the manual-approval session the clarification
specifies; the transferred archive was checksum-verified on both ends.
% src: tasks/n3_step4_record_2026-10-11.md:1021-1045; tasks/n3_step9_clarification_2026-11-15.md:20-23,88-96

\subsection{MCF checkpoints}\label{si:mcfckpt}

Table~\ref{tab:si_mcf_ckpt} lists the MCF checkpoint used for every seed of every
MCF arm. Under the step-5b set reading, a checkpoint labelled ``last'' is the
latest-epoch top-3 checkpoint, not final-epoch weights: MCF-R and MCF-A were
configured for 1,715 epochs (24,010 steps), and no checkpoint from the end of
training exists for any seed. ii-S uses the MCF-R+G checkpoints. Because MCF-R+G
chose a different checkpoint from MCF-R for seed 0, that seed was sampled
separately for MCF-R+G on TEST-B and DEV-TEST.
% src: tasks/n3_protocol.md:512-514,1530-1537; results/n3/mcf_designation_valsel.json; results/n3/n3_selection.json:285-404

\begin{table}
\caption{MCF checkpoints chosen on VALSEL. Epoch and step are as recorded by
MCF's checkpointer (epoch counted from 0). ``Decided by'' is the first criterion in
the preregistered order that separated the candidates: for MCF-R and MCF-R+G,
VALSEL picked@1, then draw@1, then any@16, then the fixed-draw VAL loss; for
MCF-A, any-of-10 at stol 1.0, then at stol 0.8, then MCF's logged validation
loss. ``Fixed-draw VAL loss'' therefore means that the candidates tied on the three
VALSEL criteria before it.}
\label{tab:si_mcf_ckpt}
\small
\begin{tabular}{@{}llllll@{}}
\toprule
arm & seed & label & epoch & step & decided by \\
\midrule
MCF-R & 0 & last & 1531 & 21,448 & fixed-draw VAL loss \\
MCF-R & 1 & last & 1563 & 21,896 & fixed-draw VAL loss \\
MCF-R & 2 & ep1067 & 1067 & 14,952 & VALSEL picked@1 \\
MCF-R+G & 0 & ep1331 & 1331 & 18,648 & VALSEL picked@1 \\
MCF-R+G & 1 & last & 1563 & 21,896 & fixed-draw VAL loss \\
MCF-R+G & 2 & ep1067 & 1067 & 14,952 & VALSEL picked@1 \\
MCF-A & 0 & ep752 & 752 & 10,542 & any-of-10 at stol 1.0 \\
MCF-A & 1 & last & 1503 & 21,056 & any-of-10 at stol 1.0 \\
MCF-A & 2 & ep1195 & 1195 & 16,744 & any-of-10 at stol 1.0 \\
\bottomrule
\end{tabular}
\end{table}
% src: results/n3/n3_selection.json:285-404 (ckpt_label, decided_by); results/n3/mcf_designation_valsel.json (runs.*.last_identical_to, epoch, global_step)
